First some notation. Let $\N = \{1,2,\ldots\}$. For $k \in \N$, define $[k] := \{1,\ldots,k\}$ and $[0] := \emptyset$. Let $I$ be the identity matrix of any size and of any vector space, which will be clear from context. For a graph $G$, let $V(G)$ denote its vertex set and $E(G)$ its edge set. For any vector $v$ with $n$ entries, let $v_i$ be its $i$th entry for all $i \in [n]$. For any matrix $A$ with $n$ rows and $m$ columns, let $A_{i,j}$ be the entry at the $i$th row and $j$th column, for all $i \in [n], j \in [m]$.

We assume the reader is familiar with the mathematical foundations of quantum computing (e.g. at the level of an upper undergraduate quantum computing class). Please see Appendix \ref{appendix:background} otherwise.